\section*{Corrected Na/Cu(001) Computational Route, Version 1.1}

\noindent\textbf{Status.}
This section supersedes the earlier Na/Cu(001) computational insert wherever
there is a conflict.  It is frozen before corrected downstream surface results.
The completed 120-SCF Cu bulk matrix is retained, but its original automatic
handoff is invalid because the selector omitted the registered energy criterion.
No adsorption, NEB, saddle, prefactor, or model-rate result is considered
executed until its hashed PASS artifact exists.

\subsection*{Distinct mathematical objects}

The ARC efficiency maximum, a reaction-path saddle, and a kinetic turnover are
not interchangeable.  The present route computes the electronic maximum along
a converged minimum-energy path.  It does not reconstruct the ARC efficiency
maximum.  A turnover analysis opens only if independently varied friction or
damping and independent rate evidence become available.

\subsection*{Authoritative files}

\begin{verbatim}
na_cu001_ci/method_protocol_v0.2.json
na_cu001_ci/validation_selection_protocol_v0.1.json
na_cu001_ci/bulk_runner_v2.py
na_cu001_ci/slab_runner_v2.py
na_cu001_ci/na_pseudo_probe_v2.py
na_cu001_ci/closure_engine_v2.py
na_cu001_ci/validate_integration_chain_v2.py
na_cu001_ci/tier_linter_v2.py
na_cu001_ci/workflow_contract_linter_v2.py
na_cu001_ci/run_computational_stage_v2.sh
.github/workflows/na-cu001-computational-route-v2.yml
\end{verbatim}

The workflow records the exact Git commit and a SHA-256 source manifest.  The
manifest is verified before integration readiness is generated.

\subsection*{Immutable inputs}

The method protocol records expected SHA-256 values for Quantum ESPRESSO 7.6,
the SSSP PBE Efficiency v2 archive, and the selected Cu and Na UPFs.  The Na
probe must locate exactly one Na UPF and read the authoritative Na cutoff values
from the registered \texttt{cutoffs.json} JSON pointer.  Frozen values are used
as assertions.  They are not substituted for missing or contradictory metadata.

\subsection*{Bulk holdout and joint gate}

Reuse the retained 20 cutoff/mesh EOS summaries from the completed 120-SCF
matrix.  Execute an additional six-point EOS at 80 Ry, 240 Ry charge-density
cutoff, and a $16\times16\times16$ k mesh over the same lattice grid.  For every
candidate compute
\[
 \Delta a_0=|a_{0,i}-a_{0,\mathrm{holdout}}|,
 \qquad
 \Delta E_0=|E_{0,i}-E_{0,\mathrm{holdout}}|.
\]
A candidate passes only if
\[
 \Delta a_0\leq0.005~\mathrm{\AA}
 \quad\text{and}\quad
 \Delta E_0\leq0.001~\mathrm{eV/atom}.
\]
Select the lexicographically smallest registered passing pair.  If none passes,
stop.  The handoff must hash-link the full bulk result.

\subsection*{Clean Cu(001) slab convergence}

Use primitive surface vectors
\[
 \mathbf A_1=(a_0/2,a_0/2,0),\qquad
 \mathbf A_2=(-a_0/2,a_0/2,0),
\]
with area $a_0^2/2$, one atom per layer, and alternating fcc stacking.  Run the
registered $4\times4\times4=64$ layer--vacuum--k cases using
\begin{verbatim}
assume_isolated = 'esm'
esm_bc = 'bc1'
\end{verbatim}
for every case.  Raw total energies of unequal atom counts are never compared.
For each vacuum and k mesh fit
\[
 E_{\mathrm{slab}}(L)=\mu_{\mathrm{slab}}L+2\epsilon_{\mathrm{surface}}.
\]
Use the frozen 1 meV per surface atom domination gate.  Retain at least seven
layers downstream even if five layers pass, so that the expanded one-sided
mobility model preserves a fixed lower surface.

\subsection*{Electrostatic consistency audit}

At the downstream layer count, repeat the selected ESM slab, repeat it at the
next larger vacuum, and run one periodic diagnostic at the selected vacuum.
PASS requires the ESM next-vacuum difference to remain within 1 meV per surface
atom.  The periodic-vs-ESM difference is reported but cannot select or retune
the route because ESM bc1 is already the common slab and downstream convention.

\subsection*{Clean relaxation and Na reference}

Relax the symmetric clean slab under the registered clean-surface z-only masks.
Freeze the final coordinates and run a separate SCF.  Require force convergence,
mirror-pair agreement, SCF convergence, and energy reproduction within 0.001 eV.

Resolve the Na UPF and its cutoffs from the pinned SSSP metadata.  Use the
componentwise maximum of the corrected Cu and metadata-derived Na cutoffs.  Run
a spin-polarized isolated Na reference in a 20 \AA\ cell.  Any archive, UPF,
metadata, unit, or hash disagreement gives HOLD.

\subsection*{Adsorption and substrate mobility}

Use one Na in a $4\times4$ primitive surface cell, nominally 0.0625 ML.  Run
hollow, bridge, and top starts at 2.0, 2.5, and 3.0 \AA\ under both models:

\begin{description}
\item[Primary:] top three Cu layers full $xyz$, one z-only buffer layer, all
remaining Cu fixed, Na full $xyz$.
\item[Expanded:] top four Cu layers full $xyz$, one z-only buffer layer, all
remaining Cu fixed, Na full $xyz$.
\end{description}

If either model has a non-hollow global minimum, write
\texttt{MECHANISM\_REVISION\_REQUIRED}, retain the landscape, and stop before
constructing the registered endpoints.

\subsection*{Endpoints and paths}

For both mobility models, verify symmetry-related hollow endpoints by relaxation,
independent SCFs, force convergence, energy agreement within 0.001 eV, site
classification, and primitive translation geometry.

Run ordinary NEB at 5, 7, and 9 images.  Select the smallest image count within
0.005 eV of every larger registered count.  Initialize CI-NEB from the hashed
selected ordinary path frames, not a fresh endpoint interpolation.  Require an
internal maximum and maximum path error no larger than 0.03 eV/\AA.

Compare the primary and expanded full-path CI barriers.  PASS requires a
difference no larger than 0.005 eV.  Otherwise preregister a larger mobile region
and do not admit a barrier.

\subsection*{Mass-weighted active-region Hessians}

At the selected endpoint and CI saddle, calculate central finite-difference
Hessians at 0.02 and 0.04 \AA\ for Na only, Na plus four nearest top-layer Cu
atoms, and Na plus those four atoms and four nearest second-layer Cu atoms.
Mass-weight the dynamical matrix before diagonalization.  Require a minimum with
no negative modes, a saddle with exactly one negative mode, finite-displacement
prefactor convergence within 20 percent, Na+4Cu to Na+8Cu convergence within
20 percent at both displacements, unstable-mode alignment, and maximum active
saddle force no larger than 0.03 eV/\AA.

Retain the Na-only and active-region prefactors and report their relative change.
Do not assume cancellation of omitted substrate modes.

\subsection*{Three-dimensional downhill connectivity}

Displace the active region in both signs of the unstable mode and relax.  A basin
match requires all of the following: periodic three-dimensional Na distance,
adsorption-height agreement, active-region RMSD including Cu, endpoint energy
agreement, site identity, and force convergence.  The two signs must reach
distinct endpoint basins.  An atom with endpoint-like $x,y$ coordinates but a
desorbed $z$ coordinate cannot pass.

\subsection*{Barrier sensitivity and rate tiers}

Do not sum image range, CI shift, and endpoint asymmetry into a statistical
uncertainty.  Retain them as diagnostics.  Define the nonprobabilistic numerical
sensitivity envelope as the maximum absolute change from the primary fixed-
geometry barrier among higher cutoff, higher density cutoff, denser k mesh,
larger vacuum, and the expanded full-path mobility result.  Set probability
coverage to null.

The final output must carry separate barrier, saddle, prefactor, rate, and
experimental tiers.  The label \texttt{COMPUTATIONAL\_FULL} is forbidden.  Five
temperature points form one nested rate curve with one independence unit; they
are not five independent observations.

\subsection*{Artifact validation and non-circular integration}

Stages 1--18 must exist as schema-correct PASS artifacts.  Each stage must
hash-link every direct dependency named in the artifact plan.  Every retained
raw file must match the raw path, size, and SHA-256 manifest.  Frozen protocol
files and the source manifest must be present.  Only after these checks pass is
Stage 19, \texttt{INTEGRATION\_READINESS.json}, generated.  The generated file
is not an input to its own validation.

\subsection*{Pre-launch verification commands}

\begin{verbatim}
python3 -m py_compile na_cu001_ci/*.py
bash -n na_cu001_ci/run_computational_stage_v2.sh
python3 na_cu001_ci/workflow_contract_linter_v2.py \
  .github/workflows/na-cu001-computational-route-v2.yml
python3 na_cu001_ci/artifact_contract_linter_v2.py \
  --plan na_cu001_ci/integration_closure_plan_v0.2.json \
  --stage-script na_cu001_ci/run_computational_stage_v2.sh
python3 na_cu001_ci/test_closure_engine.py
python3 na_cu001_ci/test_closure_engine_v2.py
python3 na_cu001_ci/test_negative_gates_v2.py
\end{verbatim}

The corrected local suite contains 23 tests, including synthetic full-chain PASS,
raw and dependency corruption, missing frozen protocols, unlinked dependencies,
desorption, unexpected minima, Cu mobility, mass weighting, tier rejection, and
workflow artifact-path and DAG checks.

\subsection*{Validation cohort and turnover boundary}

Na/Cu(001) is a development pilot and is excluded from the independent
validation cohort.  Before selecting System 2, freeze the eligible candidate
universe, exclusions, candidate-list hash, random seed, and random or stratified
sampling procedure.  Draw the system before inspecting whether its literature
appears favorable.  Retain failures and contradictions.  Any engine change after
sampling creates a new engine version and requires rerunning affected systems.

No turnover claim is permitted without an independently varied friction or
damping coordinate and independent rates.  Absence of such evidence is an
acceptable terminal result.
